\documentclass[aspectratio=169,10pt]{beamer}

% ============================================================
% CONFIGURACIÓN ROBUSTA
% ============================================================
\usepackage[utf8]{inputenc}
\usepackage[T1]{fontenc}
\usepackage[spanish,es-tabla]{babel}
\usepackage{lmodern}
\usepackage{microtype}
\usepackage{amsmath, amssymb, amsthm}
\usepackage{graphicx}
\usepackage{booktabs}
\usepackage{array}
\usepackage{multicol}
\usepackage{hyperref}
\usepackage{tikz}
\usetikzlibrary{shapes.geometric, arrows.meta, positioning, shadows, calc}
\usepackage{pifont}

% Gráficas numéricas
\usepackage{pgfplots}
\pgfplotsset{compat=1.18}

% Código
\usepackage{listings}
\usepackage{xcolor}

% ============================================================
% PALETA DE COLORES
% ============================================================
\definecolor{aerospace_blue}{RGB}{13,71,161}
\definecolor{aerospace_orange}{RGB}{230,81,0}
\definecolor{aerospace_cyan}{RGB}{0,151,167}
\definecolor{aerospace_gray}{RGB}{55,71,79}
\definecolor{code_bg}{RGB}{250,250,252}
\definecolor{code_keyword}{RGB}{13,71,161}
\definecolor{code_string}{RGB}{194,24,91}
\definecolor{code_comment}{RGB}{96,125,139}
\definecolor{highlight_yellow}{RGB}{255,243,224}

% ============================================================
% TEMA BEAMER
% ============================================================
\usetheme{Madrid}
\usecolortheme[named=aerospace_blue]{structure}

\setbeamercolor{palette primary}{bg=aerospace_blue,fg=white}
\setbeamercolor{palette secondary}{bg=aerospace_cyan,fg=white}
\setbeamercolor{palette tertiary}{bg=aerospace_orange,fg=white}
\setbeamercolor{frametitle}{bg=aerospace_blue,fg=white}
\setbeamercolor{block title}{bg=aerospace_cyan,fg=white}
\setbeamercolor{block body}{bg=highlight_yellow,fg=black}
\setbeamercolor{block title alerted}{bg=aerospace_orange,fg=white}
\setbeamercolor{block body alerted}{bg=aerospace_orange!10,fg=black}

\setbeamertemplate{navigation symbols}{}
\setbeamertemplate{page number in head/foot}[framenumber]

% ============================================================
% LOGO Y MARCA DE AGUA (CORREGIDO)
% ============================================================
\newcommand{\logofile}{logo_unii.png}

\setbeamertemplate{background}{
  \begin{tikzpicture}[remember picture,overlay]
    \node at (current page.center)
      [opacity=0.05] {\includegraphics[width=0.4\paperwidth]{\logofile}};
    \node[anchor=north east,opacity=0.7]
      at (current page.north east)
      {\includegraphics[height=0.8cm]{\logofile}};
  \end{tikzpicture}
}

% ============================================================
% LISTINGS
% ============================================================
\lstdefinestyle{pythonstyle}{
  language=Python,
  basicstyle=\ttfamily\footnotesize,
  keywordstyle=\color{code_keyword}\bfseries,
  stringstyle=\color{code_string},
  commentstyle=\color{code_comment}\itshape,
  backgroundcolor=\color{code_bg},
  frame=single,
  rulecolor=\color{aerospace_gray!30},
  frameround=tttt,
  breaklines=true,
  showstringspaces=false,
  numbers=left,
  numberstyle=\tiny\color{aerospace_gray},
  numbersep=8pt,
  xleftmargin=15pt
}

% ============================================================
% COMANDOS PERSONALIZADOS
% ============================================================
\newcommand{\importantbox}[2][Importante]{%
\begin{alertblock}{#1}#2\end{alertblock}}

\newcommand{\objectivebox}[1]{%
\begin{block}{Objetivo}#1\end{block}}

\newcommand{\dingcheck}{\textcolor{aerospace_cyan}{\ding{51}}}
\newcommand{\dingx}{\textcolor{aerospace_orange}{\ding{55}}}

% ============================================================
% FOOTER
% ============================================================
\setbeamertemplate{footline}{
  \leavevmode
  \hbox{
    \begin{beamercolorbox}[wd=.33\paperwidth,ht=2.5ex,dp=1ex,left]{author in head/foot}
      \hspace*{2ex}\insertshortauthor
    \end{beamercolorbox}
    \begin{beamercolorbox}[wd=.34\paperwidth,ht=2.5ex,dp=1ex,center]{title in head/foot}
      \insertshorttitle
    \end{beamercolorbox}
    \begin{beamercolorbox}[wd=.33\paperwidth,ht=2.5ex,dp=1ex,right]{date in head/foot}
      \insertframenumber{} / \inserttotalframenumber\hspace*{2ex}
    \end{beamercolorbox}
  }
}

% ============================================================
% DATOS DEL CURSO (CLASE 4)
% ============================================================
\title[Modelización Aeroespacial]{Modelización en Ingeniería Aeroespacial}
\subtitle{Clase 4: Estabilidad, Equilibrios y Respuesta Dinámica}
\author[Ing. Jair Molina]{Ing. Jair Molina Arce}
\institute[UNII]{Universidad Internacional de Innovación}
\date{\today}

% ============================================================
% TRANSICIÓN DE SECCIÓN
% ============================================================
\AtBeginSection[]{
\begin{frame}
\vfill
\centering
\begin{beamercolorbox}[sep=8pt,center,shadow=true,rounded=true]{title}
\usebeamerfont{title}\insertsectionhead
\end{beamercolorbox}
\vfill
\end{frame}
}

% ============================================================
% DOCUMENTO
% ============================================================
\begin{document}

% ------------------------------------------------------------
\begin{frame}
\titlepage
\end{frame}

% ------------------------------------------------------------
\begin{frame}{Agenda}
\small
\begin{itemize}
\item ¿Qué significa \textit{estabilidad} en sistemas dinámicos?
\item Puntos de equilibrio y linealización local
\item Autovalores: lectura física (modo, decaimiento, oscilación)
\item Respuesta temporal: libre vs forzada
\item Ejemplos: caída con drag y modelo $v$--$\gamma$
\item Validación e interpretación física (criterio ingenieril)
\end{itemize}
\end{frame}

% ============================================================
\section{Motivación: Simular no es entender}

\begin{frame}{De la simulación al entendimiento}
\small

\objectivebox{
Pasar de ``tengo una gráfica'' a ``sé explicar qué hace el sistema y por qué''.
}

\vspace{2mm}

\begin{block}{Contexto (Clases 1--3)}
\begin{itemize}
\item \textbf{Clase 1:} modelo físico $\rightarrow$ ecuaciones $\rightarrow$ código.
\item \textbf{Clase 2:} forma general $\dot{x}=f(x,u)$ y linealización.
\item \textbf{Clase 3:} integración numérica (Euler / RK4 / RK45).
\end{itemize}
\end{block}

\importantbox[Idea clave]{
Hoy analizamos \textbf{comportamiento} del sistema: equilibrio, estabilidad y respuesta.
\begin{itemize}
\item ¿Se regresa al equilibrio o se aleja?
\item ¿Oscila o se amortigua?
\item ¿Qué parámetro cambia la dinámica?
\end{itemize}
}

\end{frame}

% ============================================================
\section{Equilibrios y estabilidad}

\begin{frame}{Punto de equilibrio: definición}
\small

\begin{block}{Sistema no lineal}
\[
\dot{x}=f(x,u,t)
\]
\end{block}

\begin{block}{Equilibrio (para entrada constante $u=u^\star$)}
Un estado $x^\star$ es un equilibrio si:
\[
f(x^\star,u^\star,t)=0 \quad \Rightarrow \quad \dot{x}=0
\]
\end{block}

\vspace{2mm}
\textbf{Interpretación física:}
\begin{itemize}
\item Es un ``punto de reposo'' (o régimen estacionario) del sistema.
\item Si te colocas ahí y no cambias $u$, el sistema \textbf{se queda ahí}.
\end{itemize}

\importantbox[Ojo]{
En sistemas aeroespaciales, los ``equilibrios'' suelen depender de $u$:
\textbf{thrust, ángulo de ataque, empuje, control de actitud, etc.}
}

\end{frame}

\begin{frame}{Estabilidad: idea intuitiva}
\small

\begin{columns}[T]
\column{0.52\textwidth}
\begin{block}{Definición informal}
Sea $x^\star$ equilibrio. Si iniciamos cerca:
\[
x(0)=x^\star + \delta x
\]
\end{block}

\begin{itemize}
\item \textbf{Estable:} se mantiene cerca.
\item \textbf{Asintóticamente estable:} regresa al equilibrio.
\item \textbf{Inestable:} se aleja.
\end{itemize}

\column{0.48\textwidth}
\begin{center}
\begin{tikzpicture}[scale=0.95]
\draw[->,thick] (0,0) -- (4.5,0) node[right] {$t$};
\draw[->,thick] (0,0) -- (0,3.2) node[above] {$\|x-x^\star\|$};

% estable
\draw[aerospace_cyan, thick, domain=0:4, samples=60] plot(\x,{1.8+0.1*sin(deg(2*\x))});
\node[aerospace_cyan, font=\scriptsize] at (3.2,2.05) {estable};

% asintotica
\draw[exact_color, thick, domain=0:4, samples=60] plot(\x,{2.2*exp(-0.9*\x)});
\node[exact_color, font=\scriptsize] at (3.2,0.35) {asintótica};

% inestable
\draw[aerospace_orange, thick, domain=0:4, samples=60] plot(\x,{0.4*exp(0.6*\x)});
\node[aerospace_orange, font=\scriptsize] at (3.35,2.6) {inestable};

\end{tikzpicture}
\end{center}
\end{columns}

\importantbox[Traducción práctica]{
La estabilidad determina si una simulación ``explota'' por física real o por mala elección de parámetros/linealización.
}

\end{frame}

% ============================================================
\section{Linealización local y autovalores}

\begin{frame}{Linealización local alrededor del equilibrio}
\small

\begin{block}{Sistema}
\[
\dot{x}=f(x,u)
\]
\end{block}

\textbf{Expansión local alrededor de $(x^\star,u^\star)$:}
\[
x = x^\star + \delta x,\qquad u=u^\star+\delta u
\]

\begin{block}{Modelo lineal (primer orden)}
\[
\delta \dot{x} \approx A\,\delta x + B\,\delta u
\]
donde
\[
A=\left.\frac{\partial f}{\partial x}\right|_{(x^\star,u^\star)},\qquad
B=\left.\frac{\partial f}{\partial u}\right|_{(x^\star,u^\star)}
\]
\end{block}

\importantbox[Clave]{
\textbf{$A$ captura la dinámica interna local.} Para estabilidad local con $\delta u=0$ basta analizar $\delta \dot{x}=A\delta x$.
}

\end{frame}

\begin{frame}{Autovalores de $A$: lectura física}
\small

\begin{block}{Sistema lineal homogéneo}
\[
\delta \dot{x} = A\,\delta x
\]
\end{block}

\textbf{Si $A$ tiene autovalores $\lambda_i$:}
\begin{itemize}
\item $\Re(\lambda_i) < 0$ \quad $\Rightarrow$ modo se disipa \dingcheck
\item $\Re(\lambda_i) > 0$ \quad $\Rightarrow$ modo crece (diverge) \dingx
\item $\Im(\lambda_i)\neq 0$ \quad $\Rightarrow$ modo oscila (frecuencia $\sim \Im(\lambda)$)
\end{itemize}

\vspace{2mm}

\begin{columns}[T]
\column{0.58\textwidth}
\importantbox[Conexión rápida]{
\begin{itemize}
\item \textbf{Tiempo característico:} $\tau \approx 1/|\Re(\lambda)|$
\item \textbf{Frecuencia:} $\omega \approx |\Im(\lambda)|$
\end{itemize}
}

\column{0.42\textwidth}
\begin{center}
\begin{tikzpicture}[scale=0.9]
\draw[->,thick] (-2.4,0) -- (2.4,0) node[right] {$\Re(\lambda)$};
\draw[->,thick] (0,-2.2) -- (0,2.2) node[above] {$\Im(\lambda)$};

% stable points
\fill[exact_color] (-1.2, 0.0) circle (2pt);
\fill[exact_color] (-0.8, 1.2) circle (2pt);
\fill[exact_color] (-0.8,-1.2) circle (2pt);
\node[exact_color, font=\scriptsize] at (-1.35, 1.8) {estable};

% unstable point
\fill[aerospace_orange] (1.0, 0.8) circle (2pt);
\node[aerospace_orange, font=\scriptsize] at (1.25, 1.35) {inestable};

\draw[dashed, aerospace_gray] (0,-2.2) -- (0,2.2);
\node[aerospace_gray, font=\scriptsize] at (0.9, -1.9) {$\Re(\lambda)=0$};

\end{tikzpicture}
\end{center}
\end{columns}

\end{frame}

% ============================================================
\section{Respuesta temporal}

\begin{frame}{Respuesta libre vs respuesta forzada}
\small

\begin{block}{Modelo lineal}
\[
\delta \dot{x} = A\,\delta x + B\,\delta u
\]
\end{block}

\begin{columns}[T]
\column{0.5\textwidth}
\begin{block}{Respuesta libre ($\delta u=0$)}
\[
\delta \dot{x}=A\delta x
\]
Depende de:
\begin{itemize}
\item Condición inicial
\item Autovalores de $A$
\end{itemize}
\end{block}

\column{0.5\textwidth}
\begin{block}{Respuesta forzada ($\delta u\neq 0$)}
Entrada puede ser:
\begin{itemize}
\item escalón, rampa, seno
\item perturbación (viento, arrastre)
\item mando (thrust / control)
\end{itemize}
\end{block}
\end{columns}

\importantbox[Mensaje]{
En aeroespacial, muchas ``fallas'' son en realidad:
\textbf{respuesta libre inestable + perturbación pequeña}.
}

\end{frame}

\begin{frame}{Métricas prácticas de respuesta}
\small

\begin{block}{Qué sí reporta un ingeniero}
Para una salida $y(t)$ (o estado de interés):
\begin{itemize}
\item \textbf{Tiempo de establecimiento} $T_s$ (entra y permanece en banda)
\item \textbf{Sobreimpulso} $M_p$ (si hay oscilación)
\item \textbf{Amortiguamiento} (rápido/lento)
\item \textbf{Valor final} (¿tiende a algo físico?)
\end{itemize}
\end{block}

\begin{columns}[T]
\column{0.55\textwidth}
\importantbox[Regla de oro]{
Si el sistema no conserva sentido físico:
\begin{itemize}
\item revisa equilibrio y signos del modelo,
\item revisa unidades,
\item revisa parámetros,
\item luego revisa el método numérico.
\end{itemize}
}

\column{0.45\textwidth}
\begin{center}
\begin{tikzpicture}[scale=0.95]
\draw[->,thick] (0,0) -- (5.0,0) node[right] {$t$};
\draw[->,thick] (0,0) -- (0,3.2) node[above] {$y(t)$};

% response
\draw[exact_color, thick, domain=0:4.6, samples=120]
  plot(\x, {2.4*(1-exp(-1.2*\x)) + 0.35*exp(-0.7*\x)*sin(deg(3.2*\x))});

% band
\draw[dashed, aerospace_gray] (0,2.4*0.95) -- (5,2.4*0.95);
\draw[dashed, aerospace_gray] (0,2.4*1.05) -- (5,2.4*1.05);
\node[aerospace_gray, font=\scriptsize] at (4.2,2.65) {$\pm 5\%$};

% Ts marker
\draw[dashed, aerospace_cyan] (2.6,0) -- (2.6,3.0);
\node[aerospace_cyan, font=\scriptsize] at (2.6,-0.25) {$T_s$};

\end{tikzpicture}
\end{center}
\end{columns}

\end{frame}

% ============================================================
\section{Ejemplo 1: Caída con drag (relectura de estabilidad)}

\begin{frame}{Ejemplo 1: caída con drag y equilibrio}
\small

\begin{block}{Modelo (mismo caso de Clase 3)}
\[
\dot{h}=v,\qquad \dot{v}=-g-k\,v|v|
\]
\end{block}

\textbf{Interpretación:}
\begin{itemize}
\item La altura $h$ no tiene equilibrio absoluto (depende de referencia).
\item La velocidad sí tiene un régimen: \textbf{velocidad terminal}.
\end{itemize}

\begin{block}{Equilibrio en velocidad (descenso estable)}
Para caída hacia abajo ($v<0$), $|v|=-v$:
\[
\dot{v}=-g-k\,v(-v)=-g+k v^2
\]
Equilibrio cuando $\dot{v}=0$:
\[
-g+k v^2=0 \Rightarrow v^\star = -\sqrt{\frac{g}{k}}
\]
\end{block}

\importantbox[Lectura física]{
El sistema ``busca'' $v^\star$ y se estabiliza: es un equilibrio \textbf{atractor} para la velocidad.
}

\end{frame}

\begin{frame}{Estabilidad local de $v^\star$}
\small

\textbf{Dinámica 1D en velocidad (para $v<0$):}
\[
\dot{v}=f(v)=-g+k v^2
\]
Linealizamos:
\[
\delta \dot{v} \approx f'(v^\star)\,\delta v
\]
Derivada:
\[
f'(v)=2kv
\]
Evaluada en $v^\star=-\sqrt{g/k}$:
\[
f'(v^\star)=2k\left(-\sqrt{\frac{g}{k}}\right)=-2\sqrt{gk}<0
\]

\importantbox[Conclusión]{
Como $f'(v^\star)<0$, el equilibrio de velocidad es \textbf{asintóticamente estable}.
La velocidad converge a $-\sqrt{g/k}$.
}

\vspace{2mm}
\textbf{Tiempo característico aproximado:}
\[
\tau \approx \frac{1}{|f'(v^\star)|}=\frac{1}{2\sqrt{gk}}
\]

\end{frame}

% ============================================================
\section{Ejemplo 2: Modelo $v$--$\gamma$ (dinámica aero)}

\begin{frame}{Ejemplo 2: dinámica $v$--$\gamma$ (Clase 2)}
\small

\begin{block}{Modelo}
\[
\dot{v}=u - g\sin(\gamma) - k v|v|
\]
\[
\dot{\gamma}= -\frac{g}{v}\cos(\gamma)
\quad(\text{con protección si } v\approx 0)
\]
\end{block}

\textbf{Interpretación:}
\begin{itemize}
\item $v$: magnitud de velocidad (energía cinética)
\item $\gamma$: ángulo de trayectoria (cinemática del vuelo)
\item $u$: empuje/entrada equivalente (control)
\end{itemize}

\importantbox[Objetivo del análisis]{
Dado $(v^\star,\gamma^\star,u^\star)$:
\begin{itemize}
\item construimos $A=\partial f/\partial x$,
\item calculamos autovalores,
\item interpretamos estabilidad local.
\end{itemize}
}

\end{frame}

\begin{frame}[fragile]{Python: Jacobiano numérico y autovalores}
\small

\begin{lstlisting}[style=pythonstyle]
import numpy as np

def f_vgam(t, x, u, g=9.81, k=0.0020, v_eps=1e-3):
    v, gam = x
    v_safe = v if abs(v) > v_eps else (np.sign(v) * v_eps if v != 0 else v_eps)

    dv = u - g*np.sin(gam) - k*v*abs(v)
    dgam = -(g / v_safe) * np.cos(gam)
    return np.array([dv, dgam])

def jacobian_num(f, x, u, eps=1e-6):
    n = len(x)
    A = np.zeros((n, n))
    fx = f(0.0, x, u)
    for i in range(n):
        dx = np.zeros(n)
        dx[i] = eps
        A[:, i] = (f(0.0, x + dx, u) - fx) / eps
    return A

# --- punto de operación (ejemplo) ---
x_star = np.array([60.0, 0.0])   # [v*, gamma*]
u_star = 7.2                     # u*
A = jacobian_num(f_vgam, x_star, u_star)

eigA = np.linalg.eigvals(A)
print("A=\n", A)
print("eig(A)=", eigA)
\end{lstlisting}

\importantbox[Cómo leer el resultado]{
\begin{itemize}
\item $\Re(\lambda)<0$ estable local; $\Re(\lambda)>0$ inestable.
\item Si hay parte imaginaria: modo oscilatorio (espiral).
\end{itemize}
}

\end{frame}

\begin{frame}{Interpretación: ¿qué significa inestable aquí?}
\small

\begin{block}{Ejemplo de lectura}
Si un autovalor sale con $\Re(\lambda)>0$:
\begin{itemize}
\item una perturbación pequeña en $v$ o $\gamma$ \textbf{crece},
\item el modelo predice separación rápida del régimen nominal,
\item en vuelo: esto puede representar pérdida de control del ángulo de trayectoria.
\end{itemize}
\end{block}

\begin{columns}[T]
\column{0.55\textwidth}
\importantbox[Ingeniería]{
Antes de decir ``el sistema es inestable'':
\begin{itemize}
\item revisa si el punto $(x^\star,u^\star)$ realmente es un equilibrio,
\item revisa si $v^\star$ y $\gamma^\star$ son físicamente consistentes,
\item revisa unidades (rad vs grados).
\end{itemize}
}

\column{0.45\textwidth}
\begin{block}{Diagnóstico típico}
\begin{itemize}
\item $\gamma^\star$ grande $\Rightarrow$ término $g\sin\gamma$ domina.
\item $v^\star$ pequeño $\Rightarrow$ $g/v$ amplifica $\dot{\gamma}$.
\item Protecciones numéricas (eps) pueden alterar localmente.
\end{itemize}
\end{block}
\end{columns}

\end{frame}

% ============================================================
\section{Validación, sentido físico y buenas prácticas}

\begin{frame}{Checklist de validación (modo industria)}
\small

\begin{block}{Antes de confiar en una gráfica...}
\begin{enumerate}
\item \textbf{Unidades:} SI coherente (m, s, rad).
\item \textbf{Signos:} gravedad y drag con dirección correcta.
\item \textbf{Magnitudes:} parámetros dentro de rango físico.
\item \textbf{Equilibrio:} ¿de verdad $f(x^\star,u^\star)=0$?
\item \textbf{Sensibilidad:} cambia $\Delta t$ y observa convergencia.
\item \textbf{Comparación:} método simple vs método robusto (RK4/RK45).
\end{enumerate}
\end{block}

\importantbox[Regla]{
Si una conclusión depende de \textbf{un solo} $\Delta t$ o \textbf{un solo} método, aún no es conclusión.
}

\end{frame}

% ============================================================
\section{Resumen}

\begin{frame}{Lecciones clave de la Clase 4}
\small

\begin{enumerate}
\item Un \textbf{equilibrio} cumple $f(x^\star,u^\star)=0$.
\item La estabilidad local se analiza con la \textbf{linealización}:
\[
\delta \dot{x} \approx A\delta x
\]
\item Los \textbf{autovalores} de $A$ dictan el comportamiento:
\[
\Re(\lambda)\ \text{(crece/decae)},\quad \Im(\lambda)\ \text{(oscila)}
\]
\item Respuesta \textbf{libre vs forzada}: condiciones iniciales vs entradas.
\item Validación física: el criterio ingenieril es obligatorio.
\end{enumerate}

\vspace{4mm}

\begin{center}
\Large \textcolor{aerospace_blue}{\textbf{Siguiente clase:}}\\
\large Discretización de modelos y estabilidad numérica (stiffness, regiones de estabilidad)
\end{center}

\end{frame}

% ============================================================
\section*{Preguntas}

\begin{frame}
\vfill
\centering
\begin{beamercolorbox}[sep=8pt,center,shadow=true,rounded=true]{title}
\usebeamerfont{title}\Huge ¿Preguntas?
\end{beamercolorbox}
\vfill

\begin{center}
\large
\texttt{jair.molina@unii.edu.mx}\\[2mm]
\small
Material disponible en el repositorio del curso
\end{center}

\end{frame}

\end{document}
